\chapter{Borne de Cramér-Rao}

L'inégalité de Cramér-Rao fournit une minoration du risque quadratique de tout estimateur sans biais, dans un modèle paramétrique dominé satisfaisant quelques hypothèses de régularité. Elle exprime ainsi une limite intrinsèque au problème d'estimation : aucun estimateur sans biais ne peut faire mieux que la borne, quelle que soit la méthode employée pour le construire.

\textit{Ce résultat ne concerne que les estimateurs sans biais.}

\medskip
Dans toute cette section, on suppose les hypothèses de régularité suivantes, qui définissent ce que l'on appellera un \textbf{modèle régulier}:

\begin{itemize}
  \item[$(\mathcal{H}_0)$] Il existe une mesure positive $\nu$ telle que, pour tout $\theta \in \Theta$, $P_\theta$ admet une densité $f_\theta$ par rapport à $\nu$.
  \item[$(\mathcal{H}_1)$] $\Theta \subset \mathbb{R}$ est un ouvert et, pour tout $x$ de l'espace des observations $\mathcal{X}$, l'application $\theta \longmapsto f_\theta(x)$ est dérivable sur $\Theta$.
  \item[$(\mathcal{H}_2)$] Les supports des densités $f_\theta$, $\theta \in \Theta$, ne dépendent pas de $\theta$ ; on note leur valeur commune $S = \{x \in \mathcal{X} : f_\theta(x) > 0\}$.
  \item[$(\mathcal{H}_3)$] Pour tout $\theta \in \Theta$, la variable aléatoire $\frac{\partial}{\partial\theta} \ln f_\theta(X)$ est centrée sous $P_\theta$.
  \item[$(\mathcal{H}_4)$] Pour tout $\theta \in \Theta$, la variable aléatoire $\frac{\partial}{\partial\theta} \ln f_\theta(X)$ appartient à $L^2(P_\theta)$.
\end{itemize}

\begin{remarque}
L'hypothèse $(\mathcal{H}_3)$ est en réalité une conséquence d'une interversion entre
dérivation et intégration. En effet, comme $\int_S f_\theta \, d\nu = 1$ pour tout
$\theta$, si l'on peut dériver sous le signe intégral on obtient

$$
  \int_S \frac{\partial}{\partial\theta} f_\theta(x)\, d\nu(x) = 0,
  \qquad\text{soit}\qquad
  E_\theta\!\left[ \frac{\partial}{\partial\theta} \ln f_\theta(X) \right]
  = \int_S \frac{\partial_\theta f_\theta}{f_\theta}\, f_\theta \, d\nu = 0.
$$

Dans les modèles usuels, cette interversion est légitime dès que le théorème de dérivation des intégrales à paramètre s'applique (existence d'une domination intégrable de $\partial_\theta f_\theta$ au voisinage de $\theta$).
\end{remarque}

\begin{remarque}
L'hypothèse (H2) est essentielle : c'est elle qui exclut les modèles où le support
dépend du paramètre. La famille exponentielle
$\{f_\theta(x) = \theta\, e^{-\theta x}\,\mathbb{1}_{x \geq 0},\ \theta > 0\}$ a un
support $S = [0, +\infty[$ indépendant de $\theta$ et constitue un modèle régulier.
En revanche, la famille des lois uniformes $\mathcal{U}([0,\theta])$, de densité
$f_\theta(x) = \frac{1}{\theta}\mathbb{1}_{[0,\theta]}(x)$, a un support
$[0,\theta]$ qui dépend de $\theta$ : (H2) n'est pas vérifiée et la borne de
Cramér-Rao ne s'applique pas. C'est d'ailleurs pour ce type de modèle que l'on peut
obtenir des estimateurs convergeant à vitesse $1/n$, plus rapide que la vitesse
usuelle $1/\sqrt{n}$.
\end{remarque}

% ---------------------------------------------------------------------
\begin{mdframed}
\begin{definition}[Information de Fisher]
Dans un modèle paramétrique régulier, l'\textbf{information de Fisher} est la fonction $I : \Theta \to [0, +\infty[$ définie par

$$
  I(\theta)
  = E_\theta\!\left[ \left( \frac{\partial}{\partial\theta} \ln L(\theta) \right)^{2} \right]
  = \operatorname{Var}_\theta\!\left( \frac{\partial}{\partial\theta} \ln L(\theta) \right),
$$

où $L(\theta)$ désigne la vraisemblance. La variable aléatoire
$\dfrac{\partial}{\partial\theta} \ln L(\theta)$ est appelée \textbf{fonction de
score}.
\end{definition}
\end{mdframed}

L'égalité entre les deux expressions de $I(\theta)$ provient directement de (H3) :
le score étant centré, son espérance du carré coïncide avec sa variance.

\begin{remarque}[Cas multidimensionnel]
La définition est donnée pour $\theta \in \mathbb{R}$. Lorsque $\theta$ est un
vecteur de dimension $p > 1$, l'information de Fisher est la matrice $p \times p$
\[
  I(\theta)
  = \left( E_\theta\!\left[ \frac{\partial}{\partial\theta_i} \ln L(\theta)\,
    \frac{\partial}{\partial\theta_j} \ln L(\theta) \right] \right)_{1 \leq i,j \leq p},
\]
c'est-à-dire la matrice de covariance du vecteur score.
\end{remarque}

\begin{proposition}[Expression par la dérivée seconde]
On suppose de plus que, pour tout $x \in S$, l'application
$\theta \mapsto f_\theta(x)$ est de classe $C^2$ sur $\Theta$, et que l'on peut
dériver deux fois sous le signe intégral l'identité $\int_S f_\theta\, d\nu = 1$
(domination intégrable des dérivées première et seconde au voisinage de $\theta$).
Alors
\[
  I(\theta) = -\, E_\theta\!\left[ \frac{\partial^{2}}{\partial\theta^{2}} \ln L(\theta) \right].
\]
\end{proposition}

\begin{proof}
En dérivant deux fois la relation $\int_S f_\theta\, d\nu = 1$, on obtient
$\int_S \partial_\theta^2 f_\theta \, d\nu = 0$. Par ailleurs, partout sur $S$,
\[
  \frac{\partial^{2}}{\partial\theta^{2}} \ln f_\theta
  = \frac{\partial_\theta^{2} f_\theta}{f_\theta}
    - \left( \frac{\partial_\theta f_\theta}{f_\theta} \right)^{2}
  = \frac{\partial_\theta^{2} f_\theta}{f_\theta}
    - \left( \frac{\partial}{\partial\theta} \ln f_\theta \right)^{2}.
\]
En intégrant contre $f_\theta$ sous $P_\theta$ :
\[
  E_\theta\!\left[ \frac{\partial^{2}}{\partial\theta^{2}} \ln f_\theta(X) \right]
  = \int_S \partial_\theta^{2} f_\theta \, d\nu
    - E_\theta\!\left[ \left( \frac{\partial}{\partial\theta} \ln f_\theta(X) \right)^{2} \right]
  = 0 - I(\theta) = -\,I(\theta). \qedhere
\]
\end{proof}

\begin{remarque}
On ne peut pas prendre cette dernière formule comme \emph{définition} de
l'information de Fisher : elle n'est valable que sous les hypothèses de classe
$C^2$ et d'interversion ci-dessus. En revanche, la définition par la variance du
score ne requiert que (H1)--(H4). En pratique, lorsque ces hypothèses sont
satisfaites, la formule par la dérivée seconde est souvent la plus commode pour le
calcul.
\end{remarque}

\begin{remarque}[Interprétation]
L'information de Fisher mesure la quantité d'information que l'observation apporte
sur le paramètre $\theta$. La formule $I(\theta) = -E_\theta[\partial_\theta^{2}
\ln L(\theta)]$ l'identifie à la \emph{courbure moyenne} de la log-vraisemblance au
point $\theta$ : plus la log-vraisemblance est «~pointue~» autour de la vraie
valeur, plus $I(\theta)$ est grande, et plus $\theta$ est facile à estimer (la
borne de Cramér-Rao $\frac{1}{n I(\theta)}$ est alors petite). À l'inverse, une
log-vraisemblance plate, $I(\theta)$ proche de $0$, traduit que des valeurs
voisines de $\theta$ produisent des lois quasi indiscernables : l'estimation est
difficile et la borne explose. C'est l'analogue quantitatif de l'identifiabilité
vue au chapitre~1.
\end{remarque}

\begin{remarque}[Effet d'un reparamétrage]
L'information de Fisher dépend du paramétrage choisi. Si l'on pose $\eta =
\varphi(\theta)$ avec $\varphi$ un $C^1$-difféomorphisme, alors, en dérivant le
score par composition, l'information dans le paramétrage $\eta$ s'écrit
\[
  I_\eta(\eta) = \frac{I_\theta(\theta)}{\bigl(\varphi'(\theta)\bigr)^{2}},
  \qquad \theta = \varphi^{-1}(\eta).
\]
On retrouve la cohérence avec la méthode delta (\S1.4) : la borne $\frac{1}{n
I_\eta(\eta)} = \frac{(\varphi'(\theta))^{2}}{n I_\theta(\theta)}$ se transforme
exactement comme la variance asymptotique sous une transformation régulière.
\end{remarque}

% ---------------------------------------------------------------------
\begin{mdframed}
\begin{proposition}[Information de Fisher d'un $n$-échantillon]
Si le modèle $\{P_\theta, \theta \in \Theta\}$ est régulier, d'information de
Fisher $I(\theta)$, alors le modèle produit
$\{P_\theta^{\otimes n}, \theta \in \Theta\}$ est régulier, d'information de Fisher
\[
  I_n(\theta) = n\, I(\theta).
\]
\end{proposition}
\end{mdframed}

\begin{proof}
Pour un $n$-échantillon $X = (X_1, \dots, X_n)$ i.i.d., la vraisemblance se
factorise : $L_n(\theta) = \prod_{i=1}^{n} f_\theta(X_i)$, de sorte que la
log-vraisemblance, et donc le score, sont des sommes :
\[
  \frac{\partial}{\partial\theta} \ln L_n(\theta)
  = \sum_{i=1}^{n} \frac{\partial}{\partial\theta} \ln f_\theta(X_i).
\]
Les variables $\frac{\partial}{\partial\theta} \ln f_\theta(X_i)$ sont
indépendantes, identiquement distribuées, centrées (par (H3)) et de variance
$I(\theta)$ chacune. La variance d'une somme de variables indépendantes étant la
somme des variances,
\[
  I_n(\theta)
  = \operatorname{Var}_\theta\!\left( \sum_{i=1}^{n}
    \frac{\partial}{\partial\theta} \ln f_\theta(X_i) \right)
  = n\, \operatorname{Var}_\theta\!\left( \frac{\partial}{\partial\theta} \ln f_\theta(X_1) \right)
  = n\, I(\theta). \qedhere
\]
\end{proof}

% ---------------------------------------------------------------------
On considère désormais une quantité d'intérêt $g(\theta)$, où $g : \Theta \to
\mathbb{R}$ est dérivable, et un estimateur $T = T(X)$ \emph{sans biais} de
$g(\theta)$. On ajoute l'hypothèse de régularité suivante.

\begin{itemize}
  \item[\textbf{(H5)}] Pour tout $\theta \in \Theta$, on peut dériver sous le signe
        intégral :
        \[
          \frac{\partial}{\partial\theta} \int_S T(x)\, f_\theta(x)\, d\nu(x)
          = \int_S T(x)\, \frac{\partial}{\partial\theta} f_\theta(x)\, d\nu(x).
        \]
\end{itemize}

Cette écriture vaut aussi bien pour les modèles discrets que pour les modèles
dominés par la mesure de Lebesgue.

\begin{mdframed}
\begin{proposition}[Inégalité de Cramér-Rao]
Sous les hypothèses (H1)--(H4) d'un modèle régulier, si $I(\theta) > 0$ pour tout
$\theta \in \Theta$, alors tout estimateur $T_n = T(X_1, \dots, X_n)$ sans biais du
paramètre réel $g(\theta)$, vérifiant (H5) et de moment d'ordre $2$ fini
($E_\theta[T_n^2] < +\infty$), satisfait
\[
  \forall\, \theta \in \Theta, \qquad
  E_\theta\!\left[ (T_n - g(\theta))^{2} \right]
  = \operatorname{Var}_\theta(T_n)
  \geq \frac{\bigl(g'(\theta)\bigr)^{2}}{I_n(\theta)}
  = \frac{\bigl(g'(\theta)\bigr)^{2}}{n\, I(\theta)}.
\]
Le membre de droite est appelé \textbf{borne de Cramér-Rao}.
\end{proposition}
\end{mdframed}

\begin{proof}
On raisonne au niveau de l'observation $X$ (le passage au $n$-échantillon se fait
ensuite en remplaçant $I(\theta)$ par $I_n(\theta) = n I(\theta)$). Comme $T$ est
sans biais, $E_\theta[T(X)] = g(\theta)$, soit
$\int_S T(x) f_\theta(x)\, d\nu(x) = g(\theta)$. En dérivant en $\theta$ et en
utilisant (H5),
\[
  g'(\theta)
  = \int_S T(x)\, \frac{\partial}{\partial\theta} f_\theta(x)\, d\nu(x)
  = \int_S T(x)\, \frac{\partial_\theta f_\theta(x)}{f_\theta(x)}\, f_\theta(x)\, d\nu(x)
  = E_\theta\!\left[ T(X)\, \frac{\partial}{\partial\theta} \ln f_\theta(X) \right].
\]
Le score étant centré (H3), on peut soustraire $g(\theta)$ sans rien changer :
\[
  g'(\theta)
  = E_\theta\!\left[ \bigl(T(X) - g(\theta)\bigr)\, \frac{\partial}{\partial\theta} \ln f_\theta(X) \right].
\]
On reconnaît le produit scalaire $\langle h_1, h_2 \rangle_\theta =
E_\theta[h_1(X) h_2(X)]$ entre les deux variables centrées $T(X) - g(\theta)$ et
$\frac{\partial}{\partial\theta} \ln f_\theta(X)$. L'inégalité de Cauchy-Schwarz
donne
\[
  \bigl(g'(\theta)\bigr)^{2}
  \leq E_\theta\!\left[ \bigl(T(X) - g(\theta)\bigr)^{2} \right]
       \cdot E_\theta\!\left[ \left( \frac{\partial}{\partial\theta} \ln f_\theta(X) \right)^{2} \right]
  = \operatorname{Var}_\theta(T)\cdot I(\theta),
\]
la dernière égalité résultant du caractère sans biais de $T$ et de la définition de
$I(\theta)$. Comme $I(\theta) > 0$, on en déduit
\[
  \operatorname{Var}_\theta(T) \geq \frac{\bigl(g'(\theta)\bigr)^{2}}{I(\theta)}.
\]
Pour un $n$-échantillon, le même raisonnement appliqué à $L_n$ remplace $I(\theta)$
par $I_n(\theta) = n I(\theta)$, ce qui donne la borne annoncée. \qedhere
\end{proof}

\begin{remarque}[Version matricielle]
Lorsque $\theta \in \mathbb{R}^{p}$, l'inégalité devient une inégalité entre
matrices : pour tout estimateur sans biais $T$ de $\theta$,
\[
  \operatorname{Var}_\theta(T) \;\succeq\; I_n(\theta)^{-1},
\]
au sens où $\operatorname{Var}_\theta(T) - I_n(\theta)^{-1}$ est symétrique positive
(ordre de Loewner), $I_n(\theta)$ étant la matrice d'information de Fisher. En
particulier, en regardant les termes diagonaux, la variance de chaque composante
$T_j$ est minorée par le $j$-ème terme diagonal de $I_n(\theta)^{-1}$.
\end{remarque}

% ---------------------------------------------------------------------
\begin{mdframed}
\begin{definition}[Estimateur efficace]
Un estimateur $T$ sans biais de $g(\theta)$ qui réalise l'égalité dans l'inégalité
de Cramér-Rao, c'est-à-dire tel que
\[
  \operatorname{Var}_\theta(T) = \frac{\bigl(g'(\theta)\bigr)^{2}}{n\, I(\theta)}
  \qquad \forall\, \theta \in \Theta,
\]
est dit \textbf{efficace} ; on dit aussi qu'il \emph{atteint la borne de
Cramér-Rao}.
\end{definition}
\end{mdframed}

\begin{remarque}
Un estimateur efficace est, par construction, de risque quadratique minimal parmi
les estimateurs sans biais : c'est en particulier un estimateur UVMB (au sens de la
Remarque sur les estimateurs sans biais du chapitre~1). La réciproque est fausse :
un estimateur peut être UVMB sans atteindre la borne de Cramér-Rao, car cette borne
n'est pas toujours atteignable.
\end{remarque}

\begin{remarque}[Cas d'égalité]
L'égalité dans Cauchy-Schwarz a lieu si et seulement si les deux variables sont
$P_\theta$-presque sûrement proportionnelles. Un estimateur sans biais efficace
n'existe donc que si la fonction de score s'écrit
\[
  \frac{\partial}{\partial\theta} \ln L(\theta)
  = A(\theta)\,\bigl(T(X) - g(\theta)\bigr)
  \qquad P_\theta\text{-p.s.,}
\]
pour une certaine fonction $A(\theta)$. Cette contrainte n'est satisfaite que par
les modèles de type \emph{famille exponentielle}, ce qui explique que l'efficacité
soit une propriété rare.
\end{remarque}

% =====================================================================
\subsection*{Exemples}

\begin{exemple}[Modèle de Bernoulli]
Soit $(X_1, \dots, X_n)$ un $n$-échantillon i.i.d. de loi $\mathcal{B}(\theta)$,
$\theta \in {]0,1[}$, et $g(\theta) = \theta$. La densité (par rapport à la mesure
de comptage) est $f_\theta(x) = \theta^{x}(1-\theta)^{1-x}$, $x \in \{0,1\}$, d'où
\[
  \ln f_\theta(x) = x \ln\theta + (1-x)\ln(1-\theta),
  \qquad
  \frac{\partial}{\partial\theta} \ln f_\theta(x)
  = \frac{x}{\theta} - \frac{1-x}{1-\theta}
  = \frac{x - \theta}{\theta(1-\theta)}.
\]
Comme $\operatorname{Var}_\theta(X_1) = \theta(1-\theta)$,
\[
  I(\theta)
  = \operatorname{Var}_\theta\!\left( \frac{X_1 - \theta}{\theta(1-\theta)} \right)
  = \frac{\theta(1-\theta)}{\theta^{2}(1-\theta)^{2}}
  = \frac{1}{\theta(1-\theta)}.
\]
La borne de Cramér-Rao pour un estimateur sans biais de $\theta$ vaut donc
$\dfrac{1}{n I(\theta)} = \dfrac{\theta(1-\theta)}{n}$. Or la moyenne empirique
$\overline{X}_n$ est sans biais et $\operatorname{Var}_\theta(\overline{X}_n) =
\frac{\theta(1-\theta)}{n}$ : elle \textbf{atteint la borne} et est donc efficace.
\end{exemple}

\begin{exemple}[Modèle de Poisson]
Soit $(X_1, \dots, X_n)$ i.i.d. de loi $\mathcal{P}(\theta)$, $\theta > 0$, et
$g(\theta) = \theta$. On a $P_\theta(X_1 = k) = e^{-\theta}\theta^{k}/k!$, d'où
\[
  \ln f_\theta(k) = -\theta + k\ln\theta - \ln(k!),
  \qquad
  \frac{\partial}{\partial\theta} \ln f_\theta(k) = -1 + \frac{k}{\theta}
  = \frac{k - \theta}{\theta}.
\]
Comme $\operatorname{Var}_\theta(X_1) = \theta$,
\[
  I(\theta) = \frac{\operatorname{Var}_\theta(X_1)}{\theta^{2}} = \frac{1}{\theta},
  \qquad\text{borne de Cramér-Rao} = \frac{\theta}{n}.
\]
La moyenne empirique $\overline{X}_n$ est sans biais et de variance
$\frac{\theta}{n}$ : elle est efficace.
\end{exemple}

\begin{exemple}[Modèle exponentiel]
Soit $(X_1, \dots, X_n)$ i.i.d. de loi $\mathcal{E}(\theta)$, de densité
$f_\theta(x) = \theta e^{-\theta x}\mathbb{1}_{x \geq 0}$, $\theta > 0$. Alors
\[
  \ln f_\theta(x) = \ln\theta - \theta x,
  \qquad
  \frac{\partial}{\partial\theta} \ln f_\theta(x) = \frac{1}{\theta} - x,
  \qquad
  I(\theta) = \operatorname{Var}_\theta(X_1) = \frac{1}{\theta^{2}}.
\]
On ne peut pas estimer $\theta$ lui-même sans biais (cf. Exemple~1.4 du
chapitre~1). On estime donc plutôt $g(\theta) = \tfrac{1}{\theta} = E_\theta[X_1]$,
avec $g'(\theta) = -\tfrac{1}{\theta^{2}}$. La borne de Cramér-Rao est
\[
  \frac{\bigl(g'(\theta)\bigr)^{2}}{n I(\theta)}
  = \frac{1/\theta^{4}}{n/\theta^{2}} = \frac{1}{n\theta^{2}}.
\]
La moyenne empirique $\overline{X}_n$ est sans biais pour $\tfrac{1}{\theta}$ et de
variance $\frac{1}{n}\operatorname{Var}_\theta(X_1) = \frac{1}{n\theta^{2}}$ : elle
est efficace pour $g(\theta) = \tfrac{1}{\theta}$.
\end{exemple}

\begin{exemple}[Modèle gaussien]
Soit $(X_1, \dots, X_n)$ i.i.d. de loi $\mathcal{N}(\theta, \sigma^{2})$, avec
$\sigma^{2} > 0$ \emph{connue} et $g(\theta) = \theta$. La densité est
$f_\theta(x) = \frac{1}{\sqrt{2\pi\sigma^{2}}}\exp\!\bigl(-\frac{(x-\theta)^{2}}{2\sigma^{2}}\bigr)$,
d'où
\[
  \frac{\partial}{\partial\theta} \ln f_\theta(x) = \frac{x - \theta}{\sigma^{2}},
  \qquad
  I(\theta) = \frac{\operatorname{Var}_\theta(X_1)}{\sigma^{4}}
  = \frac{\sigma^{2}}{\sigma^{4}} = \frac{1}{\sigma^{2}}.
\]
La borne de Cramér-Rao vaut $\dfrac{\sigma^{2}}{n}$. La moyenne empirique
$\overline{X}_n$ est sans biais et de variance $\frac{\sigma^{2}}{n}$ : elle est
efficace. On retrouve ainsi que l'estimateur $\overline{X}_n$ de l'Exemple~1.3
(cas $\sigma^{2} = 1$, de risque $1/n$) n'est pas seulement raisonnable : il est
optimal au sens de Cramér-Rao.
\end{exemple}